% cat-lift.tex — the lift: a type is a vertex, an inst is an edge
% Rendered with the tikz MCP server (compile_tikz), which wraps this body in a
% standalone document (default preamble: tikz, amsmath, amssymb + common TikZ
% libraries) and writes docs/website/skills/mtron/assets/cat-lift.png.
%
% INK: white on a transparent background — the site's theme is dark (--maincolor
% #000000, panels #191C24). Edge labels are deliberately UNFILLED (no
% `fill=black` patch), so no text box shows against the page; labels are offset
% clear of the strokes instead. Panels use black!85 / black!75 so they read as
% slightly lighter than the page.

\begin{tikzpicture}[
  v/.style={draw, rounded corners=2pt, thick, minimum width=14mm, minimum height=8mm, font=\small},
  e/.style={-{Stealth[length=2.2mm]}, thick},
  el/.style={font=\scriptsize, inner sep=1pt},
  color=white
]
\node[v] (int)  at (0,0)      {\texttt{int}};
\node[v] (real) at (3.8,0)    {\texttt{real}};
\node[v] (bool) at (3.8,-2.1) {\texttt{bool}};
\node[v] (str)  at (0,-2.1)   {\texttt{str}};

\draw[e] (int) to[out=130,in=50,looseness=5] node[el,above=2pt]{plus} (int);
\draw[e] (real) to[out=130,in=50,looseness=5] node[el,above=2pt]{neg} (real);
\draw[e] (int) -- node[el,above=2pt]{gt} (bool);
\draw[e] (int) -- node[el,left=3pt]{as} (str);
\end{tikzpicture}
